\documentclass[11pt]{article}
\usepackage{color}
\usepackage{amsmath,amsthm,amssymb,multirow,paralist}
\usepackage[margin=0.8in]{geometry}
\usepackage{hyperref}
\begin{document}
\begin{center}
{\Large \textbf{COM S 4740/5740: Introduction to Machine Learning}\\Homework \#4}\\
\linethickness{1mm}\line(1,0){498}
\begin{enumerate}
\item Please put required code files and report into a
compressed file ``HW\#\_FirstName\_LastName.zip''
\item Unlimited number of submissions are
allowed on Canvas and the latest one will be graded.
\item {\color{red} No later submission is accepted.}
\item Please read and follow submission instructions. No exception
will be made to accommodate incorrectly submitted files/reports.
\item All students are required to typeset their reports using
latex. Overleaf
(\url{https://www.overleaf.com/learn/latex/Tutorials}) can be a
good start.
\end{enumerate}
\linethickness{1mm}\line(1,0){498}
\end{center}
%%%%%%%%%%%%%%%%%%%%%%%%%%%%%%%%%%%%%%%%%%%%%%%%%%%%%%%%%%%%%%%%%%%%%%%%%%%%%%%
%%%%%%%%%%%%%%%%%%%%%%%%%%%%%%%%%%%%%%%%%%%%%%%%%%%%%%%%%%%%%%%%%%%%%%%%%%%%%%%
\begin{enumerate}
\item (60 points) \textbf{Support Vector Machine for Handwritten
Digits Recognition}: You need to use the software package
scikit-learn
\href{https://scikit-learn.org/stable/modules/svm.html}{https://scikit-learn.org/stable/modules/svm.html}
to finish this assignment. We will use ``svm.SVC()'' to create a
svm model. The handwritten digits files are in the ``data''
folder: train.txt and test.txt. The starting code is in the
``code'' folder. In the data file, each row is a data example.
The first entry is the digit label (``1'' or ``5''), and the next
256 are grayscale values between -1 and 1. The 256 pixels
correspond to a $16\times16$ image. You are expected to implement
your solution based on the given codes. The only file you need to
modify is the ``solution.py'' file. You can test your solution by
running ``main.py'' file. Note that code is provided to compute a
two-dimensional feature (symmetry and average intensity) from
each digit image; that is, each digit image is represented by a
two-dimensional vector. These features along with the
corresponding labels should serve as inputs to your solution
functions.
\begin{enumerate}
\item (20 points) Complete the \textbf{svm\_with\_diff\_c()}
function. In this function, you are asked to try different values
of cost parameter c.
\item (20 points) Complete the \textbf{svm\_with\_diff\_kernel()}
function. In this function, you are asked to try different
kernels (linear, polynomial and radial basis function kernels).
\item (20 points) Summarize your observations from (a) and (b)
into a short report. In your report, please report the accuracy
result and total support vector number of each model. A briefly
analysis based on the results is also needed. For example, how
the number of support vectors changes as parameter value changes
and why.


\begin{color}{red}
\textbf{Results:}

\begin{table}[h]
\centering
\begin{tabular}{|c|c|c|}
\hline
\textbf{Cost Parameter ($C$)} & \textbf{Accuracy (\%)} & \textbf{Total Support Vectors} \\ \hline
0.01 & 89.39 & 1080 \\ \hline
0.1  & 96.23 & 414  \\ \hline
1    & 95.99 & 162  \\ \hline
2    & 95.99 & 129  \\ \hline
3    & 96.23 & 117  \\ \hline
5    & 96.23 & 104  \\ \hline
\end{tabular}
\caption{SVM Performance with Varying Cost Values (Linear Kernel)}
\end{table}

Analyzing the results, I noticed an inverse relationship between the cost parameter, $C$, and the number of support vectors. In an SVM, $C$ acts as a regularization parameter. A lower $C$ value allows for more margin violations, resulting in a wider margin and therefore more support vectors. A larger $C$ value penalizes misclassifications, leading to a narrower margin with fewer support vectors.

\begin{table}[h]
\centering
\begin{tabular}{|c|c|c|}
\hline
\textbf{Kernel Type} & \textbf{Accuracy (\%)} & \textbf{Total Support Vectors} \\ \hline
Linear     & 95.99 & 162 \\ \hline
Polynomial & 95.75 & 75  \\ \hline
RBF        & 96.23 & 90  \\ \hline
\end{tabular}
\caption{SVM Performance with Varying Kernels (Default Cost)}
\end{table}

The choice of kernel significantly affects the number of support vectors. More complex, non-linear kernels (such as Polynomial and RBF) allow for a more complex decision boundary. This allows the model to find more optimal, non-linear decision boundaries, and separate the classes using fewer support vectors compared to a standard linear boundary.


\end{color}

\end{enumerate}
\end{enumerate}
\end{document}